\scanpage{119}

\uline{Cor} Dans le cas de la prop 2, \deleted{la catégorie} on a
$G^0\simeq U^0\times\mathbb G_m$, $\epsilon=\mathrm{pr}_2$, $j(\lambda)=\deleted{\unsure{\mathrm{pr}(\lambda)^2|\mathrm{inj}(\lambda)^2}}\,\mathrm{inj}_{1}(\lambda)^2$. La structure
\deleted{motivique $1)\,2)\,3)$} de catégorie tensorielle de $\deleted{\unsure{G}}M_{\mathbb R}$ est
connue totalement par la connaissance de
$G$ et l'élémt de $H^{\unsure{2|1}}(\mathbb R,\deleted{\unsure{G}}U)$ caractéristique, qui \deleted{\unsure{puis.}}
\deleted{est} l'unique élémt non nul de $H^2(\mathbb R,U)$,
correspondant à l'unique élément non nul
de $H^2(\mathbb R,\mu_2)$. Il existe un \deleted{\unsure{élém}} objet $E$ de
$M_{\mathbb R}$ \deleted{\unsure{en}} de degré 1, rang 2, \deleted{sur lequel} tel que
la sous-catégorie tensorielle $\deleted{\unsure{M'}}M'$ qu'il engendre
soit déjà $\mathbb R$-tordue, et \deleted{\unsure{correspond} \unsure{au} \unsure{fait que}}
\deleted{\unsure{sur}} que $U^0$ ``y \unsure{opère|agisse} trivialmt''.
\par\hspace*{.55\linewidth}\rule{2em}{.3pt}

Dans le cas favorable $\unsure{\mathscr G|\mathcal G}_{\mathbb R}\ne\varnothing$, \deleted{\unsure{on a en gros}} \inserted{i.e. $U$ un tore}
l'ens $\deleted{\operatorname{Ob}\mathscr G_{\mathbb R}}$ \inserted{$H^0(\mathscr G_{\mathbb R})$ des objets à iso. près de $\mathscr G_{\mathbb R}$} est un
\deleted{\unsure{ens.} principal homogène} \inserted{torseur}
sous le grpe $H^{1}(\mathbb R,U)$, qui lui-$=$ est
isom. can à $ {}_2U(\mathbb R)\underset{\substack{\mathrm{isom}\\\mathrm{can.}}}{\simeq}(\mathbb Z/2\mathbb Z)^{\dim U}$. D'autre
part, à chaque élémt de $\operatorname{Ob}\mathscr G_{\mathbb R}$ correspond
une \deleted{structure de \unsure{positivité}} \inserted{structure polarisante} sur $G^0$,
définie par un hom $\mu_{2_{\mathbb C}}\mathrel{\mathop{\longrightarrow}\limits^{i_1}}G^0_{\mathbb C}$ tel que
$\epsilon_{\mathbb R}i_1=\mathrm{id}$, $i_{\unsure{1|2}}\bar i_1=j_{\mathbb R}$, et l'ens de ces structures est
égalmt un torseur sous $ {}_2U(\mathbb R)$. L'application

\[
H^0(\mathscr G_{\mathbb R})\longrightarrow\operatorname{Pol}(\mathscr G_{\mathbb R})
\]

est compatible avec \deleted{\unsure{le}} l'opération de $ {}_2U(\mathbb R)$, \uline{donc}
\uline{bijective}. Donc les $\mathbb R$-foncteurs fibres à \unsure{isom|iso|isom.}

% archive sheet number: 118
